\documentclass[a4paper,12pt]{article}

\input{../../Preambulo.tex}
\newcommand{\assunto}{FUNÇÕES MODULARES}
\newcommand{\atividade}{lista} % prova ou lista
\newcommand{\turma}{ELETROELETRÔNICA 1.\textordmasculine\ ANO}
\newcommand{\numProva}{A}
\newcommand{\professor}{Igor Martins Silva}
\newcommand{\bimestre}{3}
\newcommand{\ano}{2026}
\newcommand{\mesTexto}{agosto}
\newcommand{\mesNum}{08}
\newcommand{\dia}{27}
\newcommand{\dataTexto}{\dia\ de \mesTexto\ de \ano}
\newcommand{\dataNun}{\dia/\mesNum/\ano}

\newcommand{\titulo}{}

\ifdefstring{\atividade}{prova}
  {\renewcommand{\titulo}{Avaliação de Matemática}}
{
\ifdefstring{\atividade}{lista}
  {\renewcommand{\titulo}{Lista de exercícios de Matemática}}
{
\ifdefstring{\atividade}{as}
  {\renewcommand{\titulo}{Avaliação Somativa}}
{
\ifdefstring{\atividade}{ad}
  {\renewcommand{\titulo}{Avaliação Diagnóstica}}
  {\renewcommand{\titulo}{Atividade de Matemática}}
}}}

\newcommand{\emtipo}{%
  \ifdefstring{\tipo}{individual}{\tipo}{em \tipo}%
}

\edef\pathCapa{\currfiledir}

\newcommand{\subtitulobase}[3]{%
    % #1 = título da 1ª coluna
    % #2 = conteúdo da 1ª coluna
    % #3 = mostra número da prova? (vazio ou algo)
    
    \noindent
    \begin{minipage}{2.8cm}
        \includegraphics[width=2.8cm]{\pathCapa Logo_Cefet.png}
    \end{minipage}
    \hfill
    \begin{minipage}{\dimexpr\textwidth-6.2cm\relax}
        \begin{center}
        CEFET-MG -- Campus Contagem \\[3pt]
        \titulo\ -- \bimestre.\textordmasculine\ bimestre de \ano \\[3pt]
        \professor
    \end{center}
    \end{minipage}
    \hfill
    \begin{minipage}{3.2cm}
        \includegraphics[width=3.2cm]{\pathCapa Brasao_Brasil.png}
    \end{minipage}
    
    \begin{center}
        \vspace{15pt}
        \begin{tabular}{|C{210pt}|C{70pt}|C{210pt}|}
            \hline
            \textbf{#1} &
            \textbf{DATA} &
            \textbf{TURMA} \\
            \hline
            #2 & \dataNun & \turma \\
            \hline
        \end{tabular}
    \end{center}
    
    #3
    
    \vspace{15pt}
}

\newcommand{\subtitulolis}{%
  \subtitulobase{ASSUNTO}{\assunto}{}%
}

\newcommand{\subtitulopr}{%
  \subtitulobase{ASSUNTO}{\assunto}{\boxed{\numProva}}%
}

\newcommand{\subtituloas}{%
  \subtitulobase{DISCIPLINA}{MATEMÁTICA}{\boxed{\numProva}}%
}

\newcommand{\subtituload}{%
  \subtitulobase{DISCIPLINA}{MATEMÁTICA}{}%
}

\newcommand{\valorquestao}{}

\newcommand{\definevalorquestao}{%
    \ifdefstring{\atividade}{lista}{%
        % Não faz nada para lista
    }{%
        \ifdefstring{\modo}{objetivo}{%
            \renewcommand{\valorquestao}{\ValorQObj ponto}
        }{%
            \renewcommand{\valorquestao}{\ValorQDisc pontos}
        }%
    }%
}

\newenvironment{exercicioBanco}[1][]{%
    \ifdefstring{\atividade}{lista}{%
        \begin{exercicio}%
    }{%
        \begin{exercicio}[#1]%
    }%
}{%
    \end{exercicio}
}

\ifdefstring{\atividade}{lista}{
    \pagecolor{background-color}
}{
    
}



\hypersetup{
    pdfauthor={\professor},
    pdftitle={\titulo \ -- \assunto \ -- \turma},
}

\title{\titulo \ -- \assunto \ -- \turma}
\author{\professor}
\date{\data}

\graphicspath{
    {../../FuncoesModulares/}
}

%************* INÍCIO DO DOCUMENTO *************
\begin{document}
    \subtitulolis
        
    \phantomsection
    \addcontentsline{toc}{section}{Exercício 1}
    \input{../FM_Exercicio1_1}
    \noindent\rule{\linewidth}{1pt}  % linha horizontal fina
    
    \phantomsection
    \addcontentsline{toc}{section}{Exercício 2}
    %%%%%%%%%%%%%%%%%%%%%%%%%%%%%%%%%
%
%       EXERCÍCIO
%
%%%%%%%%%%%%%%%%%%%%%%%%%%%%%%%%%

\definevalorquestao

\begin{exercicioBanco}[\valorquestao]
Considere a função
%
\[
    \begin{array}{rcl}
        f:\bbR & \to & \bbR \\
        x & \mapsto &
        \left\{
        \begin{array}{cl}
            x+2, & \text{se } x \le -1, \\[2pt]
            1, & \text{se } -1 < x < 1, \\[2pt]
            -2x+3, & \text{se } x \ge 1.
        \end{array}
        \right.
    \end{array}
\]
%
Determine a soma dos zeros dessa função.

% Define as alternativas
\newcommand{\alternativas}{%
    \begin{center}
        \begin{tabularx}{\textwidth}{XXXXX}
            (a) \(-\dfrac{1}{2}\). &
            (b) \(-2\). &
            (c) \(0\). &
            (d) \(\dfrac{3}{2}\). &
            (e) \(1\).
        \end{tabularx}
    \end{center}
}

% Define a resposta correta
\newcommand{\resposta}{A}

% Lógica condicional para exibição
\ifdefstring{\atividade}{lista}{%
    \alternativas
    \vspace{0.5em}
    
    \noindent\textbf{Resposta:} letra \textbf{\resposta}.
}{%
    \ifdefstring{\modo}{objetivo}{%
        \alternativas
    }{%
        % modo = discursiva → não mostra alternativas
    }
}
\end{exercicioBanco}


    \noindent\rule{\linewidth}{1pt}  % linha horizontal fina
    
    \phantomsection
    \addcontentsline{toc}{section}{Exercício 3}
    %%%%%%%%%%%%%%%%%%%%%%%%%%%%%%%%%
%
%       EXERCÍCIO
%
%%%%%%%%%%%%%%%%%%%%%%%%%%%%%%%%%

\definevalorquestao

\begin{exercicioBanco}[\valorquestao]
Se \(a<0\) e \(b>0\), determine o valor de
\[
    \dfrac{|ab|}{|b|} - \dfrac{a^2}{|a|}
\]

% Define as alternativas
\newcommand{\alternativas}{%
    \begin{center}
        \begin{tabularx}{\textwidth}{XXXXX}
            (a) \(2\). &
            (b) \(-2\). &
            (c) \(0\). &
            (d) \(1\). &
            (e) \(-1\).
        \end{tabularx}
    \end{center}
}

% Define a resposta correta
\newcommand{\resposta}{C}

% Lógica condicional para exibição
\ifdefstring{\atividade}{lista}{%
    \alternativas
    \vspace{0.5em}
    
    \noindent\textbf{Resposta:} letra \textbf{\resposta}.
}{%
    \ifdefstring{\modo}{objetivo}{%
        \alternativas
    }{%
        % modo = discursiva → não mostra alternativas
    }
}
\end{exercicioBanco}


    \noindent\rule{\linewidth}{1pt}  % linha horizontal fina
    
    \phantomsection
    \addcontentsline{toc}{section}{Exercício 4}
    \input{../FM_Exercicio4_1}
    \noindent\rule{\linewidth}{1pt}  % linha horizontal fina
    
    \phantomsection
    \addcontentsline{toc}{section}{Exercício 5}
    %%%%%%%%%%%%%%%%%%%%%%%%%%%%%%%%%
%
%       EXERCÍCIO
%
%%%%%%%%%%%%%%%%%%%%%%%%%%%%%%%%%

\definevalorquestao

\begin{exercicioBanco}[\valorquestao]
Determine o gráfico da função \(f(x) = |x-2| + |x+1| - |x-1|\).

% Define as alternativas
\newcommand{\alternativas}{%
    \begin{center}
    \begin{tabular}{C{20pt}C{140pt}C{20pt}C{140pt}C{20pt}C{140pt}}
        (a) 
        &
        \begin{tikzpicture}[scale=0.5]
            % Eixos
            \draw[->] (-3,0) -- (3,0) node[right] {\(x\)};
            \draw[->] (0,-0.5) -- (0,3.5) node[above] {\(y\)};

            % Gráfico da função f(x) = |x| + |x+1| - |x-1|
            \draw[domain=-2.5:-1, thick, blue] plot (\x, {1});    % -1 \le x < 0
            \draw[domain=-1:1, thick, blue] plot (\x, {\x+2});    % -1 \le x < 0
            \draw[domain=1:2.5, thick, blue] plot (\x, {-\x+4});    % 0 \le x < 1

            % Marcação do vértice (x = -1)
            \draw[dashed] (-1,0) -- (-1,1) -- (0,1);
            \node[below] at (-1,0) {\small\(-1\)};
            \node[right] at (0,1) {\small\(1\)};
            
            \draw[dashed] (1,0) -- (1,3) -- (0,3);
            \node[below] at (1,0) {\small\(1\)};
            \node[left] at (0,3) {\small\(3\)};
            
            \draw[dashed] (2,0) -- (2,2) -- (0,2);
            \node[below] at (2,0) {\small\(2\)};
            \node[left] at (0,2) {\small\(2\)};
        \end{tikzpicture}
        &
        (b) 
        &
        \begin{tikzpicture}[scale=0.5]
            % Eixos
            \draw[->] (-3,0) -- (3,0) node[right] {\(x\)};
            \draw[->] (0,-0.5) -- (0,3.5) node[above] {\(y\)};

            % Gráfico da função f(x) = |x| + |x+1| - |x-1|
            \draw[domain=-2.5:-1, thick, blue] plot (\x, {-\x}); % x < -1
            \draw[domain=-1:1, thick, blue] plot (\x, {\x+2});    % -1 \le x < 0
            \draw[domain=1:2, thick, blue] plot (\x, {-\x+4});    % 0 \le x < 1
            \draw[domain=2:2.5, thick, blue] plot (\x, {\x});    % x \ge 1

            % Marcação do vértice (x = -1)
            \draw[dashed] (-1,0) -- (-1,1) -- (0,1);
            \node[below] at (-1,0) {\small\(-1\)};
            \node[right] at (0,1) {\small\(1\)};
            
            \draw[dashed] (1,0) -- (1,3) -- (0,3);
            \node[below] at (1,0) {\small\(1\)};
            \node[left] at (0,3) {\small\(3\)};
            
            \draw[dashed] (2,0) -- (2,2) -- (0,2);
            \node[below] at (2,0) {\small\(2\)};
            \node[left] at (0,2) {\small\(2\)};
        \end{tikzpicture}
        &
        (c)
        &
        \begin{tikzpicture}[scale=0.5]
            % Eixos
            \draw[->] (-3,0) -- (3,0) node[right] {\(x\)};
            \draw[->] (0,-0.5) -- (0,3.5) node[above] {\(y\)};

            % Gráfico da função f(x) = |x| + |x+1| - |x-1|
            \draw[domain=-2.5:-1, thick, blue] plot (\x, {-\x});    % -1 \le x < 0
            \draw[domain=-1:0, thick, blue] plot (\x, {\x+2});    % -1 \le x < 0
            \draw[domain=0:2, thick, blue] plot (\x, {2});    % 0 \le x < 1
            \draw[domain=2:2.5, thick, blue] plot (\x, {\x});    % x \ge 1

            % Marcação do vértice (x = -1)
            \draw[dashed] (-1,0) -- (-1,1) -- (0,1);
            \node[below] at (-1,0) {\small\(-1\)};
            \node[right] at (0,1) {\small\(1\)};
            
            \draw[dashed] (2,0) -- (2,2) -- (0,2);
            \node[below] at (2,0) {\small\(2\)};
            \node[left] at (0,2) {\small\(2\)};
        \end{tikzpicture}
        \\
        (d)
        &
        \begin{tikzpicture}[scale=0.5]
            % Eixos
            \draw[->] (-3,0) -- (3,0) node[right] {\(x\)};
            \draw[->] (0,-0.5) -- (0,3.5) node[above] {\(y\)};

            % Gráfico da função f(x) = |x| + |x+1| - |x-1|
            \draw[domain=0.5:2, thick, blue] plot (\x, {-\x+4});    % 0 \le x < 1
            \draw[domain=2:2.5, thick, blue] plot (\x, {\x});    % x \ge 1
            
            \draw[dashed] (1,0) -- (1,3) -- (0,3);
            \node[below] at (1,0) {\small\(1\)};
            \node[left] at (0,3) {\small\(3\)};
            
            \draw[dashed] (2,0) -- (2,2) -- (0,2);
            \node[below] at (2,0) {\small\(2\)};
            \node[left] at (0,2) {\small\(2\)};
        \end{tikzpicture}
        &
        (e) 
        &
        \begin{tikzpicture}[scale=0.5]
            % Eixos
            \draw[->] (-3,0) -- (2,0) node[right] {\(x\)};
            \draw[->] (0,-1.5) -- (0,4.5) node[above] {\(y\)};

            % Gráfico incorreto
            \draw[domain=-1.5:1, thick, blue] plot (\x, {1.5*\x + 1.5});
            \draw[domain=1:1.5, thick, blue] plot (\x, {3});

            % Guias
            \draw[dashed] (1,0) -- (1,3) -- (0,3);

            % Marcações
            \node[above left] at (-1,0) {\small\(-1\)};
            \node[above right] at (1,0) {\small\(1\)};
            \node[left] at (0,3) {\small\(3\)};
        \end{tikzpicture}
        &&
    \end{tabular}
\end{center}
%
}

% Define a resposta correta
\newcommand{\resposta}{B}

% Lógica condicional para exibição
\ifdefstring{\atividade}{lista}{%
    \alternativas
    \vspace{0.5em}
    
    \noindent\textbf{Resposta:} letra \textbf{\resposta}.
}{%
    \ifdefstring{\modo}{objetivo}{%
        \alternativas
    }{%
        % modo = discursiva → não mostra alternativas
    }
}
\end{exercicioBanco}


    \noindent\rule{\linewidth}{1pt}  % linha horizontal fina
    
    \phantomsection
    \addcontentsline{toc}{section}{Exercício 6}
    \input{../FM_Exercicio6_1}
    \noindent\rule{\linewidth}{1pt}  % linha horizontal fina
    
    \phantomsection
    \addcontentsline{toc}{section}{Exercício 7}
    %%%%%%%%%%%%%%%%%%%%%%%%%%%%%%%%%
%
%       EXERCÍCIO
%
%%%%%%%%%%%%%%%%%%%%%%%%%%%%%%%%%

\ifdefstring{\atividade}{prova}{%
    \ifdefstring{\modo}{objetivo}{%
        \renewcommand{\valorquestao}{\ValorQObj\ ponto}
    }{%
        \renewcommand{\valorquestao}{\ValorQDisc\ pontos}
    }%
}%

\begin{exercicioBanco}[\valorquestao]
Sabendo que o gráfico a seguir representa a função \(f(x) = -|x+1| + |x-1|\), determine o valor de \(y_{1}\cdot y_{2}\).
\begin{center}
    \begin{tikzpicture}[scale=0.5]
        % Eixos
        \draw[->] (-2.5,0) -- (2.5,0) node[right] {\(x\)};
        \draw[->] (0,-2.5) -- (0,2.5) node[above] {\(y\)};

        % Gráfico da função f(x) = |x-2| + |x+3|
        \draw[domain=-2:-1, thick, blue] plot (\x, {2}); % x < -1
        \draw[domain=-1:1, thick, blue] plot (\x, {-2*\x});    % x ≥ -1
        \draw[domain=1:2, thick, blue] plot (\x, {-2});    % x ≥ -1

        % Marcação do vértice (x = -1)
        \draw[dashed] (-1,0) -- (-1,2) -- (0,2);
        \node[below] at (-1,0) {\small\(x_{1}\)};
        \node[right] at (0,2) {\small\(y_{1}\)};
        \draw[dashed] (1,0) -- (1,-2) -- (0,-2);
        \node[above] at (1,0) {\small\(x_{2}\)};
        \node[left] at (0,-2) {\small\(y_{2}\)};
    \end{tikzpicture}
\end{center}


% Define as alternativas
\newcommand{\alternativas}{%
    \begin{center}
        \begin{tabularx}{\textwidth}{XXXXX}
            (a) \(-1\). &
            (b) \(-2\). &
            (c) \(-4\). &
            (d) \(-\dfrac{1}{2}\). &
            (e) \(-\dfrac{1}{4}\).
        \end{tabularx}
    \end{center}
}

% Define a resposta correta
\newcommand{\resposta}{C}

% Lógica condicional para exibição
\ifdefstring{\atividade}{lista}{%
    \alternativas
    \vspace{0.5em}
    
    \noindent\textbf{Resposta:} letra \textbf{\resposta}.
}{%
    \ifdefstring{\modo}{objetivo}{%
        \alternativas
    }{%
        % modo = discursiva → não mostra alternativas
    }
}
\end{exercicioBanco}


    \noindent\rule{\linewidth}{1pt}  % linha horizontal fina
    
    \phantomsection
    \addcontentsline{toc}{section}{Exercício 8}
    %%%%%%%%%%%%%%%%%%%%%%%%%%%%%%%%%
%
%       EXERCÍCIO
%
%%%%%%%%%%%%%%%%%%%%%%%%%%%%%%%%%

\definevalorquestao

\begin{exercicioBanco}[\valorquestao]
Seja \(f(x) = |2x - 5|\) uma função. Determine a soma dos valores de \(x\) para os quais a função assume o valor \(1\).

% Define as alternativas
\newcommand{\alternativas}{%
    \begin{center}
        \begin{tabularx}{\textwidth}{XXXXX}
            (a) \(2\). &
            (b) \(5\). &
            (c) \(10\). &
            (d) \(-2\). &
            (e) \(-5\).
        \end{tabularx}
    \end{center}
}

% Define a resposta correta
\newcommand{\resposta}{B}

% Lógica condicional para exibição
\ifdefstring{\atividade}{lista}{%
    \alternativas
    \vspace{0.5em}
    
    \noindent\textbf{Resposta:} letra \textbf{\resposta}.
}{%
    \ifdefstring{\modo}{objetivo}{%
        \alternativas
    }{%
        % modo = discursiva → não mostra alternativas
    }
}
\end{exercicioBanco}


    \noindent\rule{\linewidth}{1pt}  % linha horizontal fina
    
    \phantomsection
    \addcontentsline{toc}{section}{Exercício 9}
    %%%%%%%%%%%%%%%%%%%%%%%%%%%%%%%%%
%
%       EXERCÍCIO
%
%%%%%%%%%%%%%%%%%%%%%%%%%%%%%%%%%

\definevalorquestao

\begin{exercicioBanco}[\valorquestao]
A temperatura média \(T\), em graus Celsius, de uma cidade ao longo de um dia é dada por
%
\[
T(h)=30-|14-h|,
\]
%
em que \(h\) representa a hora do dia, com \(0\le h\le 23\). Determine em que horário a temperatura atinge seu valor máximo.

% Define as alternativas
\newcommand{\alternativas}{%
    \begin{center}
        \begin{tabularx}{\textwidth}{XXXXX}
            (a) \(0\). &
            (b) \(6\). &
            (c) \(9\). &
            (d) \(14\). &
            (e) \(23\).
        \end{tabularx}
    \end{center}
}

% Define a resposta correta
\newcommand{\resposta}{D}

% Lógica condicional para exibição
\ifdefstring{\atividade}{lista}{%
    \alternativas
    \vspace{0.5em}
    
    \noindent\textbf{Resposta:} letra \textbf{\resposta}.
}{%
    \ifdefstring{\modo}{objetivo}{%
        \alternativas
    }{%
        % modo = discursiva → não mostra alternativas
    }
}
\end{exercicioBanco}


    \noindent\rule{\linewidth}{1pt}  % linha horizontal fina
    
    \phantomsection
    \addcontentsline{toc}{section}{Exercício 10}
    %%%%%%%%%%%%%%%%%%%%%%%%%%%%%%%%%
%
%       EXERCÍCIO
%
%%%%%%%%%%%%%%%%%%%%%%%%%%%%%%%%%

\ifdefstring{\atividade}{prova}{%
    \ifdefstring{\modo}{objetivo}{%
        \renewcommand{\valorquestao}{\ValorQObj\ ponto}
    }{%
        \renewcommand{\valorquestao}{\ValorQDisc\ pontos}
    }%
}%

\begin{exercicioBanco}[\valorquestao]
Determine os valores reais de \(m\) para os quais a equação
\[
|x-1|+|x-4|=m
\]
possui infinitas soluções reais.

% Define as alternativas
\newcommand{\alternativas}{%
    \begin{center}
        \begin{tabularx}{\textwidth}{XXXXX}
            (a) \(0\). &
            (b) \(1\). &
            (c) \(2\). &
            (d) \(3\). &
            (e) \(4\).
        \end{tabularx}
    \end{center}
}

% Define a resposta correta
\newcommand{\resposta}{D}

% Lógica condicional para exibição
\ifdefstring{\atividade}{lista}{%
    \alternativas
    \vspace{0.5em}
    
    \noindent\textbf{Resposta:} letra \textbf{\resposta}.
}{%
    \ifdefstring{\modo}{objetivo}{%
        \alternativas
    }{%
        % modo = discursiva → não mostra alternativas
    }
}
\end{exercicioBanco}


    \noindent\rule{\linewidth}{1pt}  % linha horizontal fina
    
    \phantomsection
    \addcontentsline{toc}{section}{Exercício 11}
    %%%%%%%%%%%%%%%%%%%%%%%%%%%%%%%%%
%
%       EXERCÍCIO
%
%%%%%%%%%%%%%%%%%%%%%%%%%%%%%%%%%

\definevalorquestao

\begin{exercicioBanco}[\valorquestao]
Considere as afirmações a seguir e responda se são verdadeiras ou falsas.
\begin{enumerate}[(i)]
    \item O conjunto solução da inequação \(|x - 2| < 1\) é \(S = ]-\infty,1]\;\cup\;[3,\infty[\).
    %
    \item Para todo \(x, y \in \bbR\) vale que \(|x - y| = |y - x|\).
    %
    \item \(-|-2x| = -2|x|\), para todo \(x \in \bbR\).
\end{enumerate}
% Define as alternativas
\newcommand{\alternativas}{%
    \begin{center}
        \begin{tabularx}{\textwidth}{XXXXX}
            (a) (F, F, V). &
            (b) (F, V, V). &
            (c) (V, F, F). &
            (d) (F, V, F). &
            (e) (F, F, F).
        \end{tabularx}
    \end{center}
}

% Define a resposta correta
\newcommand{\resposta}{B}

% Lógica condicional para exibição
\ifdefstring{\atividade}{lista}{%
    \alternativas
    \vspace{0.5em}
    
    \noindent\textbf{Resposta:} letra \textbf{\resposta}.
}{%
    \ifdefstring{\modo}{objetivo}{%
        \alternativas
    }{%
        % modo = discursiva → não mostra alternativas
    }
}
\end{exercicioBanco}


    \noindent\rule{\linewidth}{1pt}  % linha horizontal fina
    
    \phantomsection
    \addcontentsline{toc}{section}{Exercício 12}
    \input{../FM_Exercicio12_1}
    \noindent\rule{\linewidth}{1pt}  % linha horizontal fina
    
    \phantomsection
    \addcontentsline{toc}{section}{Exercício 13}
    %%%%%%%%%%%%%%%%%%%%%%%%%%%%%%%%%
%
%       EXERCÍCIO
%
%%%%%%%%%%%%%%%%%%%%%%%%%%%%%%%%%

\ifdefstring{\atividade}{prova}{%
    \ifdefstring{\modo}{objetivo}{%
        \renewcommand{\valorquestao}{\ValorQObj\ ponto}
    }{%
        \renewcommand{\valorquestao}{\ValorQDisc\ pontos}
    }%
}%

\begin{exercicioBanco}[\valorquestao]
Sejam \(S_{1}\) e \(S_{2}\) os conjuntos soluções das seguintes equações modulares, respectivamente: 
\begin{multicols}{2}
    \begin{enumerate}[(i)]
        \item \(|2x+1| = x-4\),
        \item \(|x-5| = 5-x\).
    \end{enumerate}
\end{multicols}
Determine \(S_{1} \cup S_{2}\).

% Define as alternativas
\newcommand{\alternativas}{%
    \begin{center}
        \begin{tabularx}{\textwidth}{XXX}
            (a) \(\{5\}\). &
            (b) \(\{1,5\}\). &
            (c) \(\varnothing\). \\[5pt]
            (d) \(\{x \in \bbR \mid x \le 5\}\). &
            (e) \(\bbR\).
        \end{tabularx}
    \end{center}
}

% Define a resposta correta
\newcommand{\resposta}{D}

% Lógica condicional para exibição
\ifdefstring{\atividade}{lista}{%
    \alternativas
    \vspace{0.5em}
    
    \noindent\textbf{Resposta:} letra \textbf{\resposta}.
}{%
    \ifdefstring{\modo}{objetivo}{%
        \alternativas
    }{%
        % modo = discursiva → não mostra alternativas
    }
}
\end{exercicioBanco}


    \noindent\rule{\linewidth}{1pt}  % linha horizontal fina
    
    \phantomsection
    \addcontentsline{toc}{section}{Exercício 14}
    %%%%%%%%%%%%%%%%%%%%%%%%%%%%%%%%%
%
%       EXERCÍCIO
%
%%%%%%%%%%%%%%%%%%%%%%%%%%%%%%%%%

\definevalorquestao

\begin{exercicioBanco}[\valorquestao]
Determine o conjunto solução, em \(\bbR\), da inequação \(1 < |x - 1| < 3\).

% Define as alternativas
\newcommand{\alternativas}{%
    \begin{center}
        \begin{tabularx}{\textwidth}{XXX}
            (a) \([0,2]\) &
            (b) \(]-2,0[\cup]2,4[\) &
            (c) \(]-2,4[\) \\[5pt]
            (d) \(]-\infty,0[\) &
            (e) \(]-\infty,0[\cup]2,\infty[\)
        \end{tabularx}
    \end{center}
}

% Define a resposta correta
\newcommand{\resposta}{B}

% Lógica condicional para exibição
\ifdefstring{\atividade}{lista}{%
    \alternativas
    \vspace{0.5em}
    
    \noindent\textbf{Resposta:} letra \textbf{\resposta}.
}{%
    \ifdefstring{\modo}{objetivo}{%
        \alternativas
    }{%
        % modo = discursiva → não mostra alternativas
    }
}
\end{exercicioBanco}


    \noindent\rule{\linewidth}{1pt}  % linha horizontal fina
    
    \phantomsection
    \addcontentsline{toc}{section}{Exercício 15}
    %%%%%%%%%%%%%%%%%%%%%%%%%%%%%%%%%
%
%       EXERCÍCIO
%
%%%%%%%%%%%%%%%%%%%%%%%%%%%%%%%%%

\definevalorquestao

\begin{exercicioBanco}[\valorquestao]
Resolva, em \(\bbR\), a inequação \(|x - 2| + |x - 4| \ge 6\).

% Define as alternativas
\newcommand{\alternativas}{%
    \begin{center}
        \begin{tabularx}{\textwidth}{XXX}
            (a) \(\bbR\). &
            (b) \(\varnothing\). &
            (c) \(]-\infty,0]\). \\[5pt]
            (d) \([6,\infty[\). &
            (e) \(]-\infty,0] \cup [6,\infty[\).
        \end{tabularx}
    \end{center}
}

% Define a resposta correta
\newcommand{\resposta}{E}

% Lógica condicional para exibição
\ifdefstring{\atividade}{lista}{%
    \alternativas
    \vspace{0.5em}
    
    \noindent\textbf{Resposta:} letra \textbf{\resposta}.
}{%
    \ifdefstring{\modo}{objetivo}{%
        \alternativas
    }{%
        % modo = discursiva → não mostra alternativas
    }
}
\end{exercicioBanco}


    \noindent\rule{\linewidth}{1pt}  % linha horizontal fina
    
    \phantomsection
    \addcontentsline{toc}{section}{Exercício 16}
    %%%%%%%%%%%%%%%%%%%%%%%%%%%%%%%%%
%
%       EXERCÍCIO
%
%%%%%%%%%%%%%%%%%%%%%%%%%%%%%%%%%

\ifdefstring{\atividade}{prova}{%
    \ifdefstring{\modo}{objetivo}{%
        \renewcommand{\valorquestao}{\ValorQObj\ ponto}
    }{%
        \renewcommand{\valorquestao}{\ValorQDisc\ pontos}
    }%
}%

\begin{exercicioBanco}[\valorquestao]
Um pedestre caminha em uma avenida retilínea. Sua distância, em quilômetros, até uma farmácia localizada no quilômetro \(1\) é dada por
\[
|x-1|,
\]
e sua distância até um hospital localizado no quilômetro \(6\) é dada por
\[
|x-6|.
\]

Sabe-se que a soma dessas distâncias é igual a \(5,5\) quilômetros. Determine as posições possíveis do pedestre.


% Define as alternativas
\newcommand{\alternativas}{%
    \begin{center}
        \begin{tabularx}{\textwidth}{XXXXX}
            (a) \(\{0,7\}\). &
            (b) \(\left\{\dfrac{3}{4},\dfrac{25}{4}\right\}\). &
            (c) \(\left\{\dfrac{11}{4},\dfrac{11}{2}\right\}\). &
            (d) \(\{0,11\}\). &
            (e) \(\{7,25\}\).
        \end{tabularx}
    \end{center}
}

% Define a resposta correta
\newcommand{\resposta}{B}

% Lógica condicional para exibição
\ifdefstring{\atividade}{lista}{%
    \alternativas
    \vspace{0.5em}
    
    \noindent\textbf{Resposta:} letra \textbf{\resposta}.
}{%
    \ifdefstring{\modo}{objetivo}{%
        \alternativas
    }{%
        % modo = discursiva → não mostra alternativas
    }
}
\end{exercicioBanco}


    \noindent\rule{\linewidth}{1pt}  % linha horizontal fina
    
    \phantomsection
    \addcontentsline{toc}{section}{Exercício 17}
    %%%%%%%%%%%%%%%%%%%%%%%%%%%%%%%%%
%
%       EXERCÍCIO
%
%%%%%%%%%%%%%%%%%%%%%%%%%%%%%%%%%

\definevalorquestao

\begin{exercicioBanco}[\valorquestao]
Uma fábrica considera aceitável que o comprimento \(x\), em milímetros, de uma peça satisfaça simultaneamente:
\begin{itemize}
    \item distância máxima de \(0{,}8\) mm do valor \(10\);
    \item distância mínima de \(0{,}3\) mm do valor \(9\).
\end{itemize}
Determine todos os valores possíveis de \(x\).

% Define as alternativas
\newcommand{\alternativas}{%
    \begin{center}
        \begin{tabularx}{\textwidth}{XXX}
            (a) \([9,3\,;\,10,8]\). &
            (b) \([9,2\,;\,8,7]\cup[9,3\,;\,10,8]\). &
            (c) \([9,2\,;\,8,7)\cup(9,3\,;\,10,8]\). \\[5pt]
            (d) \([9,2\,;\,10,8]\). &
            (e) \((-\infty,8,7]\cup[9,3,\infty)\).
        \end{tabularx}
    \end{center}
}

% Define a resposta correta
\newcommand{\resposta}{A}

% Lógica condicional para exibição
\ifdefstring{\atividade}{lista}{%
    \alternativas
    \vspace{0.5em}
    
    \noindent\textbf{Resposta:} letra \textbf{\resposta}.
}{%
    \ifdefstring{\modo}{objetivo}{%
        \alternativas
    }{%
        % modo = discursiva → não mostra alternativas
    }
}
\end{exercicioBanco}


    \noindent\rule{\linewidth}{1pt}  % linha horizontal fina
    
\end{document}
%*************** FINAL DO DOCUMENTO ***************
